\noindent\unskip\SourceLeaf{18}Also
\[\begin{aligned}B&\equiv-NN+1\equiv-(b-c)^2+1\\C&\equiv-NN+b-c\equiv-(b-c)^2+b-c\\D&\equiv-NN-(b-c)\equiv-(b-c)^2-(b-c)\end{aligned}\]
\[\begin{array}{rrrr}
b-c\equiv0&B\equiv1&C\equiv0&D\equiv0\\1&0&0&1\\2&0&1&0.
\end{array}\]
\par\noindent\hspace*{\parindent}Ist also 3 Res. $n$, so ist $N$ durch 3 theilbar.
\Needspace{10\baselineskip}
\begingroup\small
\begin{center}BEMERKUNGEN ZU DEN VORSTEHENDEN NOTIZEN ÜBER CUBISCHE UND BIQUADRATISCHE RESTE.\end{center}
\noindent\hspace*{\parindent}Die Notiz [1] ist von \textsc{Gauss} auf das Vorsatzblatt des Einbandes seines Handexemplars der Disquisitiones arithmeticae geschrieben. Die Angaben unter [2] finden sich in einem mit der Überschrift »Uraniae sacrum« versehenen Hefte, und ebenda folgen zwei Seiten später die beiden allgemeinen Sätze [3]. Die Entwicklungen unter [4] sind einigen losen Zetteln entnommen.

\textsc{Gauss} giebt zu Beginn seiner Commentatio prima über die Theorie der biquadratischen Reste an, dass er mit dem Jahre 1805 begonnen habe, die Theorie der cubischen und biquadratischen Reste zu durchforschen. Die hier unter [2] und [3] gegebenen Entwicklungen folgen in dem genannten Hefte unmittelbar auf Notizen, welche das Datum 1804 tragen. Die »Observatio venustissima« [1], welche weniger weit greifend ist, als die »Criterien« [2] und [3], kann zeitlich nicht wohl später liegen als letztere. Man wird es demnach hier mit den ältesten \textsc{Gauss}'schen Untersuchungen und damit überhaupt mit den ältesten Urkunden über die höheren Reciprocitätsgesetze zu thun haben.

Wie die Observatio so sind vermuthlich auch die Criterien Inductionsresultate. Es ist aber sehr bemerkenswerth, dass schon hier bei den Criterien die Zerlegung der quadratischen Form $aa+27bb$ in ihre beiden complexen Factoren und damit der Zahl $p$ in ihre beiden complexen Primtheiler auftritt. Es war bereits hiermit die Bahn gewonnen, auf welcher \textsc{Gauss} späterhin wenigstens das biquadratische Reciprocitätsgesetz in seine einfachste Gestalt kleidete. Bei gewissen Beweisansätzen des allgemeinen cubischen Reciprocitätssatzes (siehe die gleich folgenden Notizen), welche jedenfalls vor 1809 liegen, ist übrigens \textsc{Gauss} wenigstens am Beginn allein vom Gebrauch rationaler Zahlen ausgegangen.

Die zu Anfang von [2] unter 3) gemachte Angabe lässt sich als Specialfall aus dem ersten allgemeinen Satze unter [3] entnehmen. Letzterer liefert direct den Reciprocitätssatz in der allgemeinsten und einfachsten Gestalt, falls $p=3m+1$ und $q=3n-1$ ist. Der cubische Charakter des Primtheiler $\pi_1=\dfrac{a+3b\sqrt{-3}}{2}$ in Bezug auf $q$ ist nämlich durch
\[\left[\frac{\pi_1}{q}\right]\equiv(A-B\sqrt{-3})^{q-1}\quad\;(\mathrm{mod.}\ q)\]
